\section{共同编年：从潢河到河西的并立秩序（907--1227）}

这一章不把辽与西夏排成宋朝北、西两侧的背景。它们各有王室、军队、城镇、征收和道路；宋朝进入同一时间线，只在它实际改变这些条件的地方出现。日期的骨架主要来自《辽史》本纪、营卫志与《宋史》夏国传，并须同《续资治通鉴长编》、宋人使辽记录以及碑刻、城址和黑水城文书互校。官修史保留了政权的年序，却往往把多地行动压成皇帝的一道命令；下文能说的是各方可见的选择与其空间条件，而非替任何一方补写统一的意图。

\subsection{多中心的辽与党项边地（907--1038）}

九〇七年耶律阿保机取得汗位、九一六年称帝建国，首先要处理的是潢河流域的部众、营卫与通向东北和燕地的关系，而不只是改一个国号。九三六年辽支持石敬瑭建立后晋，十六州随之进入辽的政治结构；九三八年经营幽州为南京，使城池、农业腹地和南北道路成为多个中心之间可调度、也难以完全掌握的资源。九四七年辽军一度进入开封而后撤，恰好说明军队抵达都城不等于能在那里征粮、安置和维持秩序。《辽史》与五代材料都记录了这段进退；至于城内反应和撤军的比例，材料不足以写成整齐的单一原因。

宋朝九六〇年建国后，河东仍是辽、北汉与宋之间的缓冲。九七九年宋灭北汉并在高梁河受辽军击退，边界由此不再隔着一个政权。辽景宗、其后的萧太后与圣宗要面对的并非抽象的南方威胁，而是燕云城镇、河北道路、骑兵机动和补给的实际配合；九八六年宋军多路北伐，更使这些条件经受检验。宋、辽双方本纪都保存了胜败，兵数、伤亡与将领责任却随战后叙事而变，不能用一边的数字替另一边定案。

同一时期，党项李氏在夏州、灵州与黄河西北的城寨间经营部众。九八二年李继迁脱离宋朝羁縻，九〇年代又在受封与独立行动之间周旋；一〇〇四年继迁死后，李德明承接的是道路、草场、互市与军事网络，而不是一份已经稳定的领土。宋方册封确实是可核的外交动作，却不能将“受封”直接译成地方控制已经实现。宋人叙述在此极为丰富，黑水城的西夏文文书、河西寺院和城址则使另一部分制度与日常可见；它们的出土集中、定年和释读参差，也限定了我们所能概括的范围。

\subsection{盟约、战争与城寨成本（1004--1082）}

\label{evt:liao-xia-chanyuan}
一〇〇四年辽圣宗与萧太后南至澶州，宋真宗至前线；次年盟约落实。澶州附近的河道、渡口、粮道与河北堡寨，是双方从对峙转向使节、银绢和边市安排的现实背景。宋方所谓岁币、辽方礼制中的称谓都带有各自的政治语言，不能单独推出臣属关系。盟约所能稳定的是一套需要不断维护的通行、侦察和供给秩序，而不是把边界变成无人地带；此后辽军行动与宋方边防调整仍提醒各地居民，和平本身也有驻军和征发的成本。

李元昊在一〇三二年继位后改变王权象征，一〇三八年称帝建国。宋朝以“僭”界定此事，是宋廷外交立场，不足以代替党项国家化的过程。一〇四〇年的三川口、一〇四一年的好水川、一〇四二年的定川寨，把高原水源、山地道路、堡寨与后方粮运带到前景；一〇四四年宋夏议和，赐予与互市成为压低军事消耗的安排。史书可以确证战争和和议的次序，却不能可靠给出两军的精确人数，亦不能让宋方战报替西夏说明战略选择。

一〇八〇年宋廷向西北进取，两年后一〇八二年永乐城失守。城并非只是一枚攻守棋子：取水、运粮、驻兵与征发会随城寨的建立而延长。宋方对战败的责任归属与伤亡统计需要相互对读，西夏一侧的沉默也不能被填作毫无损耗的胜利。这里可见的制度得失是，宋朝试图以更前置的据点改变边防，代价却由沿线运输、守军和地方财力持续承担；这种压力会在此后与辽、金及西夏的关系中反复出现。

\subsection{旧秩序的分解与河西的转换（1114--1227）}

一一一四年完颜阿骨打起兵反辽，次年金建国；一一二〇年宋金结盟，一一二二年燕京失守。辽面对的并不是一个由宋朝单独制造的新敌人，而是东北军政网络、燕云防线与南方外交同时松动。一一二四年耶律大石西行，随后形成西辽；一一二五年天祚帝被俘，辽在原有核心区域的统治终结。皇室覆亡不等于契丹人、燕云城市或辽代制度经验同时消失；它们会以迁徙、降附、地方官署和西行网络的不同速度改变。西辽的后续更须同波斯文和中亚材料互校，不能只以中原王朝的亡国叙事收束。

一二〇六年铁木真获成吉思汗称号，西夏北方的力量平衡改变；一二〇九年蒙古攻夏后出现和议与联姻安排。河西走廊的绿洲、关隘、寺院与商路，使这一关系不只是草原军队对一个国号的压力。一二二六年蒙古大举入夏，次年西夏覆亡。蒙古、宋、金及后出的元代材料对战争的因果和战果并不一致，尤其没有条件把城破、人口损失或接收范围写成精确总数。较可确认的是，新的统治者接收了部分城镇和交通节点；黑水城文书所见的语言、佛教、印刷与行政实践则提示，政权终结没有使河西与宁夏的居民经验同时归零。
